% ============================================================
%  Übungsaufgaben Template (my_dhbw Stil, uebung_*.tex)
%  Header "Übungsblatt", grün eingefärbte <details>-Lösungen.
%  Renderer: pdflatex (zweimal)
% ============================================================
\documentclass[11pt, a4paper]{article}

\usepackage[ngerman]{babel}
\usepackage{fontspec}
\setmainfont[
  Path=/usr/share/fonts/TTF/,
  Extension=.ttf,
  BoldFont=DejaVuSans-Bold,
  ItalicFont=DejaVuSans-Oblique,
  BoldItalicFont=DejaVuSans-BoldOblique
]{DejaVuSans}
\setsansfont[
  Path=/usr/share/fonts/TTF/,
  Extension=.ttf,
  BoldFont=DejaVuSans-Bold,
  ItalicFont=DejaVuSans-Oblique,
  BoldItalicFont=DejaVuSans-BoldOblique
]{DejaVuSans}
\setmonofont[Path=/usr/share/fonts/TTF/,Extension=.ttf]{DejaVuSansMono}
\usepackage[top=2cm, bottom=2cm, left=2.4cm, right=2.4cm, footskip=1cm]{geometry}
\usepackage{amsmath, amssymb}
\usepackage{xcolor}
\usepackage{enumitem}
\usepackage{fancyhdr}
\usepackage{lastpage}
\usepackage{tabularx}
\usepackage{booktabs}
\usepackage{microtype}
\usepackage{parskip}

\usepackage{array}
\usepackage{longtable}
\usepackage{ltablex}
\keepXColumns
\usepackage{colortbl}
\usepackage[most,breakable]{tcolorbox}
\usepackage{listings}
\usepackage[normalem]{ulem}
\usepackage{pdflscape}
\usepackage{titlesec}
\usepackage[colorlinks=true, linkcolor=accent, urlcolor=accent]{hyperref}

\renewcommand{\familydefault}{\sfdefault}

% ── Theme ─────────────────────────────────────────────────────
\definecolor{accent}{RGB}{0, 60, 113}
\definecolor{primaryblue}{RGB}{0, 60, 113}
\definecolor{loesung}{RGB}{0, 100, 0}
\definecolor{codebg}{RGB}{245, 245, 245}
\definecolor{figurebg}{RGB}{232, 241, 255}
\definecolor{tablehintbg}{RGB}{240, 255, 240}
\definecolor{solutionbg}{RGB}{240, 252, 240}
\definecolor{rulegray}{RGB}{180, 180, 180}
\definecolor{tableheadbg}{RGB}{0, 60, 113}

% ── Header/Footer ─────────────────────────────────────────────
\pagestyle{fancy}
\fancyhf{}
\renewcommand{\headrulewidth}{0.4pt}
\fancyhead[L]{\footnotesize\color{accent}\nichttitle}
\fancyhead[R]{\footnotesize\color{accent}\nichtsubtitle}
\fancyfoot[R]{\footnotesize\color{accent}Blatt \thepage\,/\,\pageref{LastPage}}

% ── Typografie ────────────────────────────────────────────────
\titleformat{\section}
    {\color{accent}\large\bfseries}{}{0pt}{}
    [\vspace{-4pt}\color{accent}\rule{\linewidth}{1.5pt}]
\titleformat{\subsection}
    {\normalsize\bfseries\color{accent!85!black}}{}{0pt}{}{}
\titleformat{\subsubsection}
    {\small\bfseries\color{accent!70!black}}{}{0pt}{}{}
\titlespacing*{\section}{0pt}{10pt}{3pt}
\titlespacing*{\subsection}{0pt}{6pt}{2pt}
\titlespacing*{\subsubsection}{0pt}{4pt}{1pt}

\setlength{\parindent}{0pt}
\setlength{\parskip}{6pt}
\renewcommand{\arraystretch}{1.2}
\setlist[itemize]{noitemsep, topsep=3pt, parsep=0pt, leftmargin=1.4em}
\setlist[enumerate]{noitemsep, topsep=3pt, parsep=0pt, leftmargin=1.6em}

% ── Boxen: solutionbox = Lösungsbox in grün ──────────────────
\newtcolorbox{figurebox}[1]{
  colback=figurebg, colframe=accent!50,
  title={Abbildung: #1}, fonttitle=\small\bfseries,
  breakable, before skip=6pt, after skip=6pt
}
\newtcolorbox{tablehintbox}[1]{
  colback=tablehintbg, colframe=green!50!black!50,
  title={Tabelle: #1}, fonttitle=\small\bfseries,
  breakable, before skip=6pt, after skip=6pt
}
\newtcolorbox{solutionbox}[1]{
  enhanced,
  colback=solutionbg, colframe=loesung,
  title={#1}, fonttitle=\small\bfseries,
  coltitle=white, colbacktitle=loesung,
  breakable, before skip=8pt, after skip=8pt
}

\lstset{
  basicstyle=\ttfamily\small,
  backgroundcolor=\color{codebg},
  breaklines=true,
  frame=single, framerule=0pt,
  xleftmargin=4pt, xrightmargin=4pt,
  aboveskip=6pt, belowskip=6pt,
  keepspaces=true
}

\providecommand{\nichttitle}{Übungsblatt}
\providecommand{\nichtsubtitle}{}

\renewcommand{\nichttitle}{Analysis 2}
\renewcommand{\nichtsubtitle}{Übungsblatt 5}


\begin{document}

\begin{center}
{\Large\bfseries\color{accent}\nichttitle}
\ifx\nichtsubtitle\empty\else
  \\[2pt]{\normalsize\color{accent!80!black}\nichtsubtitle}
\fi
\\[2pt]
\color{accent}\rule{0.4\linewidth}{0.8pt}\color{black}
\end{center}
\vspace{4pt}

% ── BODY ──────────────────────────────────────────────────────

\section{Übungsblatt 5 — Komplexe Zahlen}


\noindent\textcolor{rulegray}{\rule{\linewidth}{0.4pt}}


\paragraph{Aufgabe 1 — Grundlagen und Konsequenzen \textit{(10 Punkte)}}\mbox{}\\[2pt]

a) Definieren Sie die imaginäre Einheit i und die Menge ℂ formal. Geben Sie für eine komplexe Zahl z = a+ib Real- und Imaginärteil an. \textit{(3 P)}


b) Begründen Sie anhand der Definition von i, warum die Gleichung x² = −25 in ℝ keine, in ℂ aber genau zwei Lösungen besitzt. Geben Sie diese an. \textit{(3 P)}


c) Zeigen Sie allgemein: Für jede komplexe Zahl z gilt z·z̄ = |z|². Folgern Sie daraus, dass z·z̄ stets reell und nicht-negativ ist. \textit{(4 P)}


\textbf{Lösung:}


a) i ist definiert durch i = √−1, d. h. i² = −1. Damit ist ℂ = \{a+ib | a,b ∈ ℝ\}. Für z = a+ib heißt Re(z) = a der Realteil und Im(z) = b der Imaginärteil.


b) In ℝ existiert keine Zahl mit Quadrat −25, da Quadrate reeller Zahlen ≥ 0 sind. In ℂ gilt mit i² = −1:

(5i)² = 25·i² = −25 und (−5i)² = 25·i² = −25.

Lösungen sind also x₁ = 5i und x₂ = −5i.


c) Sei z = a+ib, dann z̄ = a−ib. Es folgt:

z·z̄ = (a+ib)(a−ib) = a² − (ib)² = a² − i²b² = a² + b².

Da a,b ∈ ℝ, ist a²+b² ∈ ℝ und ≥ 0. Außerdem gilt nach Definition |z| = √(a²+b²), also z·z̄ = |z|². ☐



\noindent\textcolor{rulegray}{\rule{\linewidth}{0.4pt}}


\paragraph{Aufgabe 2 — Rechenoperationen \textit{(15 Punkte)}}\mbox{}\\[2pt]

Gegeben: z₁ = 3+4i, z₂ = 1−2i, z₃ = −2+2i.


a) Berechnen Sie z₁+z₂, z₁·z₂ und z₁² jeweils in kartesischer Form. \textit{(5 P)}


b) Berechnen Sie z₁/z₂ in der Form a+ib. Geben Sie den genauen Rechenweg über Erweiterung mit dem konjugiert Komplexen an. \textit{(5 P)}


c) Bestimmen Sie |z₃| sowie das Argument φ = arg(z₃) und schreiben Sie z₃ in Polarform r·$e^{iφ}$. \textit{(5 P)}


\textbf{Lösung:}


a) Addition:

z₁+z₂ = (3+1)+i(4−2) = 4+2i.


Multiplikation mit Formel (a₁a₂−b₁b₂)+i(a₁b₂+a₂b₁):

z₁·z₂ = (3·1 − 4·(−2)) + i(3·(−2) + 1·4) = (3+8) + i(−6+4) = 11 − 2i.


Quadrieren mit (a²−b²)+2ab·i:

z₁² = (9−16) + 2·3·4·i = −7 + 24i.


b) Erweitern mit z̄₂ = 1+2i:

z₁/z₂ = (3+4i)(1+2i) / ((1−2i)(1+2i))

= (3+6i+4i+8i²) / (1+4)

= (3−8 + 10i) / 5

= (−5+10i)/5

= −1 + 2i.


c) Betrag:

|z₃| = √((−2)² + 2²) = √8 = 2√2.


Argument: z₃ liegt im 2. Quadranten (Re < 0, Im > 0). arctan(2/(−2)) = arctan(−1) = −π/4, durch Korrektur um +π:

φ = π − π/4 = 3π/4.


Polarform: z₃ = 2√2 · $e^{i·3π/4}$.



\noindent\textcolor{rulegray}{\rule{\linewidth}{0.4pt}}


\paragraph{Aufgabe 3 — Eulersche Formel und Potenzieren \textit{(20 Punkte)}}\mbox{}\\[2pt]

a) Schreiben Sie z = 1+i in Polarform und bestimmen Sie damit z⁸ in kartesischer Form. \textit{(6 P)}


b) Eine Wechselgröße werde durch w(t) = 5·$e^{i·ωt}$ mit ω = π/3 (Einheit 1/s) modelliert. Berechnen Sie w(2) als komplexe Zahl in kartesischer Form, sowie |w(t)| für beliebiges t. \textit{(6 P)}


c) Begründen Sie mit der Eulerschen Formel, weshalb die Multiplikation einer komplexen Zahl z = r·$e^{iφ}$ mit i einer Drehung um 90° entspricht und der Betrag dabei unverändert bleibt. \textit{(8 P)}


\textbf{Lösung:}


a) z = 1+i: r = √(1+1) = √2, φ = arctan(1/1) = π/4 (1. Quadrant). Also z = √2 · $e^{iπ/4}$.


Mit zˣ = rˣ·$e^{iφx}$:

z⁸ = (√2)⁸ · $e^{i·8·π/4}$ = 2⁴ · $e^{i·2π}$ = 16 · (cos 2π + i sin 2π) = 16·1 = 16.


Also z⁸ = 16.


b) w(2) = 5·$e^{i·π/3·2}$ = 5·$e^{i·2π/3}$ = 5·(cos(2π/3) + i·sin(2π/3))

= 5·(−1/2 + i·√3/2) = −5/2 + i·5√3/2.


Betrag: |w(t)| = |5·$e^{iωt}$| = 5·|$e^{iωt}$| = 5·√(cos²(ωt)+sin²(ωt)) = 5·1 = 5 (zeitunabhängig).


c) i hat in Polarform den Betrag 1 und das Argument π/2, denn:

i = 0 + 1·i = 1·(cos(π/2) + i·sin(π/2)) = $e^{iπ/2}$.


Mit den Potenzregeln der Exponentialfunktion folgt:

z·i = r·$e^{iφ}$ · $e^{iπ/2}$ = r · $e^{i(φ+π/2}$).


Daraus liest man ab:

\begin{itemize}
  \item Der neue Betrag ist r — also unverändert.
  \item Das neue Argument ist φ + π/2 — also eine Drehung um genau 90° gegen den Uhrzeigersinn.
\end{itemize}

Geometrisch entspricht die Multiplikation in der Polarform einer Multiplikation der Beträge und einer Addition der Argumente; mit |i|=1 wirkt nur die Winkeladdition. ☐



\noindent\textcolor{rulegray}{\rule{\linewidth}{0.4pt}}


\paragraph{Aufgabe 4 — Transfer: n-te Einheitswurzeln \textit{(20 Punkte)}}\mbox{}\\[2pt]

Im Kapitel wurde gezeigt, dass für z = r·$e^{iφ}$ gilt: zⁿ = rⁿ·$e^{iφn}$. Wenden Sie dies an, um alle Lösungen w ∈ ℂ der Gleichung w⁴ = 16 zu bestimmen.


a) Argumentieren Sie zunächst, warum jede Lösung den Betrag |w| = 2 haben muss. \textit{(4 P)}


b) Schreiben Sie die rechte Seite 16 in Polarform und nutzen Sie aus, dass $e^{iφ}$ = $e^{i(φ+2πk}$) für k ∈ ℤ gilt, um alle möglichen Argumente φₖ der Lösungen zu finden. \textit{(8 P)}


c) Geben Sie alle vier Lösungen sowohl in Polar- als auch in kartesischer Form an. Verifizieren Sie eine davon durch Einsetzen in w⁴ = 16. \textit{(8 P)}


\textbf{Lösung:}


a) Sei w = r·$e^{iφ}$. Dann w⁴ = r⁴·$e^{i4φ}$. Aus w⁴ = 16 folgt für die Beträge:

|w⁴| = r⁴ = |16| = 16, also r = √[4]\{16\} = 2 (r ≥ 0 als Betrag).


b) 16 = 16·$e^{i·0}$ = 16·$e^{i·2πk}$ für k ∈ ℤ. Damit:

r⁴·$e^{i4φ}$ = 16·$e^{i·2πk}$ ⇒ r=2, 4φ = 2πk ⇒ φₖ = πk/2.


Für k = 0,1,2,3 erhält man vier verschiedene Werte (k=4 wiederholt k=0 mod 2π):

φ₀ = 0, φ₁ = π/2, φ₂ = π, φ₃ = 3π/2.


c) Polarform: wₖ = 2·$e^{iπk/2}$.



{\small
\begin{tabularx}{\linewidth}{XXX}
\hline
\rowcolor{tableheadbg}
{\color{white}\textbf{k}} & {\color{white}\textbf{Polarform}} & {\color{white}\textbf{kartesisch}} \\[2pt]
\hline
\endhead
\hline
\endfoot
0 & 2·$e^{i·0}$ & 2 \\
1 & 2·$e^{iπ/2}$ & 2i \\
2 & 2·$e^{iπ}$ & −2 \\
3 & 2·$e^{i·3π/2}$ & −2i \\
\end{tabularx}
}


Verifikation für w₁ = 2i: (2i)⁴ = 2⁴·i⁴ = 16·(i²)² = 16·(−1)² = 16·1 = 16. ✓



\noindent\textcolor{rulegray}{\rule{\linewidth}{0.4pt}}


\paragraph{Aufgabe 5 — Fehler finden \textit{(15 Punkte)}}\mbox{}\\[2pt]

Eine Studentin rechnet folgendes Beispiel mit a = −8+6i, b = 12−9i und c = 1+i (Werte aus dem Kapitel). Sie behauptet:


\begin{lstlisting}
Schritt 1:  c² = (1+i)² = 1² + i² = 1 − 1 = 0
Schritt 2:  |a| = √(−8² + 6²) = √(−64+36) = √−28
Schritt 3:  a/b = (−8+6i)/(12−9i) = −8/12 + 6i/(−9i) = −2/3 − 2/3
\end{lstlisting}


a) Identifizieren Sie in jedem der drei Schritte den Fehler und benennen Sie die jeweils verletzte Regel. \textit{(9 P)}


b) Berechnen Sie alle drei Ausdrücke korrekt. Bei a/b reicht das Ergebnis aus dem Kapitel — bestätigen Sie dieses durch Erweiterung mit b̄. \textit{(6 P)}


\textbf{Lösung:}


a) Schritt 1: Beim Quadrieren einer Summe muss die binomische Formel (x+y)² = x²+2xy+y² benutzt werden. Der gemischte Term 2·1·i fehlt.

Schritt 2: Es wurde −8² statt (−8)² gerechnet. Wegen der Operatorpräferenz steht −8² für −(8²) = −64; korrekt ist (−8)² = 64.

Schritt 3: Komplexe Division ist nicht komponentenweise. Es muss mit dem konjugiert Komplexen erweitert werden, denn der Nenner muss reell gemacht werden.


b) Korrekt:


c² = (1+i)² = 1² + 2·1·i + i² = 1 + 2i − 1 = 2i.


|a| = √((−8)² + 6²) = √(64+36) = √100 = 10.


a/b = (−8+6i)/(12−9i) erweitern mit b̄ = 12+9i:

Zähler: (−8+6i)(12+9i) = −96 − 72i + 72i + 54i² = −96 + 54·(−1) = −150.

Nenner: (12−9i)(12+9i) = 12² + 9² = 144+81 = 225.

Damit a/b = −150/225 = −2/3 (rein reell), in Übereinstimmung mit dem Kapitel.



\noindent\textcolor{rulegray}{\rule{\linewidth}{0.4pt}}





% ──────────────────────────────────────────────────────────────

\end{document}
